%!TEX program = uplatex
\RequirePackage{plautopatch}

\documentclass[uplatex,dvipdfmx]{jlreq}

% 余白調整（1ページ収め）
\usepackage[top=18mm,bottom=20mm,left=20mm,right=20mm]{geometry}

\usepackage{amsmath,amssymb}
\usepackage{url}

\pagestyle{empty}

% 行間少し詰める
\renewcommand{\baselinestretch}{0.96}

%--------------------------------
% 講演マクロ
%--------------------------------

\newcommand{\talk}[3]{%
\noindent
#1\quad #2\par
\hspace*{2em}#3\par\smallskip
}

%--------------------------------
% タイトル情報
%--------------------------------

\newcommand{\EwMnumber}{第37回}
\newcommand{\EwMtitle}{数学者のための分子生物学入門\\[2mm]\large 新しい数学を造ろう}
\newcommand{\EwMdate}{2006年6月30日（金）14:30～18:00 ～ 7月1日（土）10:30～17:00}
\newcommand{\EwMplace}{東京都文京区春日1-13-27 中央大学理工学部}

%--------------------------------
% タイトル表示
%--------------------------------

\newcommand{\EwMtitleblock}{

\vspace*{-5mm}

\begin{center}

{\Large ENCOUNTER with MATHEMATICS}

\vspace{2mm}

{\large 数学との遭遇}

\vspace{4mm}

{\Large \bfseries \EwMnumber\ \EwMtitle}

\vspace{4mm}

\EwMdate

\vspace{2mm}

於：\EwMplace

\end{center}

\vspace{2mm}

}

\begin{document}

\EwMtitleblock

%--------------------------------
% プログラム
%--------------------------------

\section*{プログラム}

\subsection*{6月30日（金）}

\talk
{14:30～16:00}
{分子生物学の非常に基礎的な事柄についての解説}
{加藤 毅 氏（京大・理・数学）}

\talk
{16:30～18:00}
{生体内ネットワークの数理モデル}
{阿久津 達也 氏（京大・化研・バイオインフォマティクスセンター）}

\subsection*{7月1日（土）}

\talk
{10:30～12:00}
{タンパク質の折り畳みに関する最適化問題}
{岡本 祐幸 氏（名大・理・物理）}

\talk
{13:40～15:10}
{遺伝子の系統樹と系統ネットワーク}
{斎藤 成也 氏（国立遺伝学研究所）}

\talk
{15:30～17:00}
{蛋白質間相互作用ネットワークの構造と進化}
{田中 博 氏（東京医科歯科大・疾患生命科学研究部）}

\bigskip

17:10～　懇親会（ワイン・パーティー）

%--------------------------------
% 案内文
%--------------------------------

\bigskip

別紙の趣旨に沿った集会の第37回を以上のような予定で開催いたします。  
非専門家向けに入門的な講演をお願い致しました。  
多く方々のご参加をお待ちしております。  
講演者による講演内容へのご案内を添付いたしますのでご覧下さい。

%--------------------------------
% 連絡先
%--------------------------------

\section*{連絡先}

112-8551 東京都文京区春日1-13-27 中央大学理工学部数学教室 

tel : 03-3817-1745  

ENCOUNTER with MATHEMATICS  
homepage : \url{http://www.math.chuo-u.ac.jp/ENCwMATH}  

三松 佳彦 : yoshi@math.chuo-u.ac.jp / 高倉 樹 : takakura@math.chuo-u.ac.jp

\end{document}
